\section{Questions suggested by the calculations}
\label{sec:questions}

The ideal volume proposal already has an established measure-theoretic
realization in~\eqref{eq:polyhedralpartition}, and the tree methods described
in Section~\ref{sec:hyperideal} also compute bordered genus-zero volumes.
The following fifteen questions start from the results and explicit
hypotheses of the preceding chapters. They isolate further research,
not claims that no related results exist. The local formulas, external
theorem inputs, conditional estimates, and finite computations supporting
each starting point have different roles. Resolving these questions is
not a prerequisite for publishing those results with their stated limits.

\begin{question}[Direct hyperideal integration]
For four labeled geodesic boundaries of fixed positive lengths, can the
domains and local density~\eqref{eq:hyperidealdensity} be integrated
directly, with all cell and symmetry factors, to recover the already
known volume $2\pi^2+\frac12\sum_iL_i^2$ in~\eqref{eq:basevolumes}?
Which change of variables makes the polynomiality visible in these
hyperideal coordinates? The positive-boundary extension described in
Section~\ref{sec:hyperideal} uses a selected pair $L_1<L_2$ to give a
full-measure decomposition into tree and double-tree cells
\cite[Theorem~3]{Zonneveld}. Equal-length volume values follow by
continuity. These spine/tree cells are not identified with the
hyperideal angle cells.
\end{question}

\begin{question}[A usable measure on general polyhedral cells]
The inverse-Poisson formula~\eqref{eq:pfaffiandensity} is already a local
density on a symplectic leaf. Can it be expressed in independent
dihedral-angle variables in a way that makes higher-dimensional cell
integrals useful? A calculation must specify the coordinate Jacobian,
the fixed hyperideal boundary lengths or compact cone openings, and the
geometric chamber and quotient. In the ideal case the disintegrated
measure~\eqref{eq:deltameasure} supplies a simple model; a determinant
in redundant angle coordinates does not by itself supply the replacement.
\end{question}

\begin{question}[Gluing and topology]
Can a polyhedral description be combined with Fenchel--Nielsen gluing to
give a practical calculation for arbitrary $(g,n)$? A proposed formula
must specify both the cutting multiplicity and the full mapping-class
quotient, including stabilizer weights. The twist periods in~\eqref{eq:fn}
remove the Dehn-twist ambiguity along a chosen cut; they do not make
all pants-length coordinates a fundamental domain for moduli space.
\end{question}

\begin{question}[Curvature corrections in period coordinates]
Can the correction class $\Xi_\mu$ in~\eqref{eq:Xiclass} be represented
by an explicit first curvature correction to the WP form in
the actual finite-area conformal coordinates? Equation~\eqref{eq:Frickefirstcorrection}
answers the local four-point calculation in a blown-up character chart,
whose conversion to conformal coordinates is still required.
How does the local density behave
in a joint limit toward a balanced collision wall, where the
final-chamber expansion is not uniform in the weights?
\end{question}

\begin{question}[Beyond the regular quantum degeneration]
Can Proposition~\ref{prop:regularquantumAiry}, for the pole-cleared
equal-label family and bounded real transition parameter, be extended
uniformly over compact sets of nondegenerate Euclidean tetrahedral
shapes, with a geometric description of the required scalar normalization?
Can a geometric state
integral in that regime produce the Thurston--Deligne--Mostow measure
on the boundary-shape space? The model~\eqref{eq:modelpartition}
specifies a classical comparison, not an identification with an
existing quantum invariant. The exact Liouville
identity~\eqref{eq:exactLiouvillecurvature} and the torsion
limit~\eqref{eq:torsionDMlimit} identify two precise boundary
constructions; the auxiliary-tetrahedron obstruction and the
different regimes in~\eqref{eq:areaLiouvillecharge} explain
what remains to connect them to that joint asymptotic.
The finite one-screening fusion calculation in
Section~\ref{sec:screenedfusion} computes the finite period and
residue-sector comparison, with a logarithmic limit at balanced
weights. Its direct infinite-spectral derivation uses the
continuation and tail-control input~\eqref{eq:screeningtailinput};
Section~\ref{sec:normalizedsixj} fixes the quantum-group
normalization, proves the rank-two obstruction, and computes
the classical rank-three fold in the regular compact model.
Section~\ref{sec:quantumAiry} proves the regular-family formula with
its first correction using the exact integration-coordinate dictionary,
pole clearing, contour and tail estimates. Its finite exact-integrand
evaluations test that normalization on specified inputs. They give no
uniform theorem for arbitrary shapes, unbounded transition parameter,
or global gluing. The fixed-shape root-of-unity theorems in
Section~\ref{sec:quantumdegeneration} have their own admissible label
sequences; the refined Chen--Murakami prefactor has stricter sequence
hypotheses than the leading volume limit. Neither supplies this joint
degeneration or identifies its measure.
\end{question}

\begin{question}[Uniform quantum Dehn filling]
For a knot complement admitting the positive geometric optimization
of Section~\ref{sec:dehnresponse}, can an exact quantum surgery
representation be controlled uniformly in a complex neighborhood
of complete peripheral holonomy, with a nonzero geometric contribution
and bounds on all competing terms? Proposition~\ref{prop:dehncriticalvalue}
already fixes every classical long-slope coefficient, and
\eqref{eq:dehngaussiantorsion} specifies the torsion check. The further
issue is the analytic input of
Proposition~\ref{prop:dehnquantumconditional}, especially in a joint
limit where the slope varies with the quantum level.
\end{question}


\begin{question}[Certified realization and quantitative angle bounds]
For the planar problem of Section~\ref{sec:ipolybridge}, can the
finite-centering certificate and short-step iteration estimate
\eqref{eq:ipolyStepBound} be combined with positive-feasibility search,
controlled inexact solves, and geometric error certification into a
complete complexity theorem with explicit input and precision dependence?
The existing iteration estimate starts from a specified centered
feasible point; the dense finite examples do not bound the cost of
finding that point or certify their floating-point errors.
For a three-manifold angle
polytope whose maximum is interior, which additional geometric hypotheses
bound its smallest angle in terms of the combinatorial input and a
quantitative feasible margin? The planar Laplacian sign structure does
not extend automatically to the full three-dimensional Hessian.
\end{question}

\begin{question}[A continuum comparison of second variations]
Can a family of refined polyhedral functionals, with specified geometric
convergence and normalization, have a second-variation limit related to
a WP potential? The exact cotangent identity, Schur complement, and
Gaussian determinant in Section~\ref{sec:ipolybridge} provide discrete
objects for such a comparison, with the stated boundary grounding or
quotient by constants. Control of their refinement limits and an
explicit identification of the target metric are still required;
finite determinant identities supply neither.
\end{question}

\begin{question}[Continuation of a coefficient Laplace transform]
For a fixed nonempty multicurve on the closed surface of
Section~\ref{sec:coefficientdistributions}, does $D_C(s)$ admit a
meromorphic continuation beyond the proved half-plane $\Re s>1$?
The kernel in~\eqref{eq:coeffLaplacedefinition} is $e^{-sL}$.
A proof must control coefficient extraction in a suitable topology
on the continued character-valued expression: extraction is not
bounded for the uniform norm on unitary characters. Pointwise
continuation of twisted Selberg functions does not justify that step.
\end{question}

\begin{question}[Control of the globally filtered remainder]
Proposition~\ref{prop:globalonecrossing} gives a bounded global
cancellation of every one-crossing primitive and all its iterates.
Can the signed residual $\mathcal R_X$ in~\eqref{eq:globalresidual}
be estimated with constants controlled by the geometry of $X$?
The exact obstruction~\eqref{eq:globalobstruction} rules out a
perfect selector in the specified class of fixed linear functionals.
Length-dependent cancellation or the nonlinear diagonal extraction
in~\eqref{eq:globaldiagonal} retains information that such a selector
discards. The latter is an exact length-valued construction, whose
quantitative stability under approximation remains an additional task.
\end{question}

\begin{question}[Effective selection of connected supports]
Can the signed coefficient relations produce the smoothed estimate
\eqref{eq:countsmoothhypothesis} with specified $\eta,m$ and geometric
constants? The elementary lattice remainder, the stable finite support
projection~\eqref{eq:countsupportfourier}, and the constructive isolation
estimate are available. The global one-crossing filter and the
length-valued Euler product now provide concrete aggregate objects.
An estimate for the higher-intersection residual or for the nonlinear
first-atom extraction, uniform over the growing family of supports and
strong enough to yield a counting saving, remains to be proved.
Fixed-coefficient isolation does not supply this aggregate estimate.
\end{question}

\begin{question}[From weighted curve counts to mixing]
Can the filling-pair coefficient construction of
Section~\ref{sec:pairconverse} produce a norm-controlled estimate
for the weighted orbit kernel~\eqref{eq:twopointkernel} without
assuming mixing? Proposition~\ref{prop:twopointunfolding} gives the
exact unfolding with the reciprocal polar Jacobian. Under the assumed
estimate for terminal tests $H_G=(h+V)G$, with $F\in L^2(\mu)$,
$G\in\operatorname{Dom}(V)$ and $0<\delta\le h$,
Theorem~\ref{thm:twopointfullrate} preserves
the counting exponent $\delta$. A bound for one endpoint test alone
has the weaker de-averaging conclusion of
Theorem~\ref{thm:twopointconverse}. The norm-controlled two-endpoint
estimate, its Jacobian weights, and the noncompact tail remain to be
established by this coefficient route for the principal-stratum kernel;
a scalar area or length count is insufficient.
Section~\ref{sec:genus2markedtorus}
settles the separate affine return estimate on the nonsquare $(2,2)$
stratum, and Section~\ref{sec:coefficientobservables} gives selected
primitive-arc coefficient observables with uniform geometric norms.
Can their class be enlarged to control the full two-endpoint kernel,
with quantitative stability of the nonlinear coefficient extraction
and passage from auxiliary hyperbolic lengths to area on the
simultaneous mapping-class quotient? Can the calculation be extended
beyond the marked-torus locus with its distinct invariant measure?
\end{question}

\begin{question}[The sharp norm-ball remainder]
Can the qualitative $o_X(L\log L)$ remainder of
Proposition~\ref{prop:strictconvexlittleo} be made effective through a
decay modulus satisfying~\eqref{eq:primitiveDini}, or through cancellation
in~\eqref{eq:primitiveexactnormalizedmobius}? The inverse-flatness
formula~\eqref{eq:inverseflatmicroscopic} quantifies a local lower
obstruction, while~\eqref{eq:dyadicerrorinverse} gives a test independent
of the leading coefficient. Determining whether the primitive remainder
is $O_X(L)$, or has an unbounded slowly varying factor, remains open here.
Section~\ref{sec:arithmeticgeometry} makes the two inputs explicit:
$\Delta_B(u)=O(u/\log u)$ together with any fixed Mertens power
saving gives $O_X(L)$, whereas an integrable geometric modulus and
the prime number theorem already give $o_X(L)$.
The all-point target is itself conditional on compatibility with the
local geometry: under the supporting-row hypotheses of
Proposition~\ref{prop:inverseflatrow}, it requires
$W(h/L^2)=O(1/\log L)$ at the contact scales $L=n/h$.
The double-exponential model in Section~\ref{subsec:flatnessmodels}
violates that necessary condition. No such supporting configuration
or compatible modulus is established for the two computed tori.
The maintained length-$8192$ campaign evaluates the finite diagnostic
\eqref{eq:allPointLogDiagnostic} in hyperbolic length units; the
unit-shortest-vector gauge requires rescaling by the systole.
It reproduces the displayed counts and block values, but the original
precision and sensitivity archive remains unrecovered. Its new
sampled precision checks and conditional sensitivity envelopes must
retain their own provenance. Neither a finite range nor an assumed
arithmetic saving determines the missing geometric modulus.
\end{question}

\begin{question}[Beyond closure operations]
Which positive constant series have a geometric partition or a useful
functional equation behind them? The figure-eight kernel
\eqref{eq:universalkernel} and the twist kernel~\eqref{eq:Kkernel} provide
two concrete models. Their different index sets suggest separating
global curve identities from identities for a distinguished twist orbit.
\end{question}

\begin{question}[A boundary phase-space formulation]
Can the endpoint-and-direction matching relations in
Section~\ref{sec:matching} be organized into an integration scheme for
general filling curves? Orthogeodesic lengths give the canonical local
geometry, while the endpoint phase variables retain the missing gluing
information. A useful result would specify the invariant measure,
matching multiplicities and convergence of both the itinerary sum and
the phase-variable integrals. A cusp version must also retain the
truncation convention of Section~\ref{sec:cusps}; its
almost-everywhere cusp-normal-ray discussion assumes no geodesic
boundary. Surfaces with geodesic boundary require the additional
exiting-ray contributions.
\end{question}

These questions concern concrete integrals or pointwise identities. The
distinctions established above---fixed parameters, geometric chambers,
correct curve types, and specified measures---are part of their statements,
rather than technicalities to be added after a formal calculation.

\paragraph{Computational support.}
The maintained normalization, Hessian/barrier, quantum/surgery,
regular-family exact-integrand, counting/mixing, and finite
quantum/geodesic companions test specified identities or parameter sets.
The last companion is described in
\path{experiments/quantum_geodesics/README.md}, with result path
\path{artifacts/quantum_geodesics/result.json}. Its local trace, twist,
cusp and quantum formulas retain the finite scope of the named groups.
The torus campaign adds
finite length and event data. These computations are mapped to the
questions in \path{claims/section32-appendixb.json}, with their result
paths and limits. They do not implement direct general-cell integration,
global quantum gluing, a stable infinite coefficient extractor, or an
independent proof of a mixing or lattice-remainder rate. Such extensions
require their stated analytic or geometric inputs before a numerical
experiment could test a well-defined instance.
